\documentclass[11pt]{article}
\usepackage[a4paper,margin=2.2cm]{geometry}
\usepackage{xcolor}
\usepackage{fontspec}
\usepackage{xeCJK}
\usepackage{enumitem}
\setmainfont{Times New Roman}
\setCJKmainfont{SimSun}
\setlist[itemize]{leftmargin=1.4em,itemsep=0.35em,topsep=0.35em}
\setlength{\parindent}{0pt}
\setlength{\parskip}{0.55em}
\pagestyle{plain}
\begin{document}
{\color[HTML]{1F5FD2}\LARGE\bfseries 注意力缺陷多动障碍（ADHD）课堂管理\par}
{\color[HTML]{667085}\small 星轨原创教育资源：用于家庭与学校共同制定支持计划。\par}
\vspace{0.8em}
{\color[HTML]{111827}\large\bfseries 课堂环境调整\par}
\begin{itemize}
\item ADHD 学生常在持续注意、冲动控制、组织材料和时间管理方面遇到困难，需要外部结构支持。
\item 座位可靠近教师和任务材料，远离门窗与高流量通道，桌面只保留当前任务所需物品。
\end{itemize}
{\color[HTML]{111827}\large\bfseries 任务设计\par}
\begin{itemize}
\item 把长作业拆成五到十分钟的小段，每完成一段就用勾选、简短反馈或计划内微休息作为结束信号。
\item 使用计时器、检查清单、书面步骤和颜色标记，让学生知道先做什么、完成到哪里、作业交到哪里。
\end{itemize}
{\color[HTML]{111827}\large\bfseries 行为支持\par}
\begin{itemize}
\item 对期望行为进行即时且具体的强化，例如开始动笔、举手发言、按清单整理材料。
\item 提前安排合理活动量，如分发材料、站立书写、伸展或短距离跑腿，避免把所有移动都视为问题行为。
\end{itemize}
\vfill
{\color[HTML]{667085}\footnotesize 本材料为应用内原创教育内容，不能替代医疗、法律或正式评估建议。\par}
\end{document}
